% v6 section. Numbers are generated \vsix macros (E6, E6b evidence); the per-surface
% table is generated (v6-e6-table.tex). Conferral is a design, not a result.
\section{Where the Guarantee Ends}\label{sec:v6boundary}

The guarantee in Sections~\ref{sec:v6stop}--\ref{sec:v6cost} has three boundaries. It holds only at a sink that can
refuse a stale epoch. It covers only what the specification names. It says nothing about how much authority people
choose to grant. We measure the first, the second was measured above, and the third we state as an open problem with a
predeclared design.

\subsection{Which surfaces can be fenced}

Our sink is cooperative by construction: it is our code, and it checks the epoch because we wrote it to. A deployment's
effects land elsewhere, in databases, coordination stores, file systems, object stores and orchestration APIs. For each
such surface the question is whether \emph{the surface itself} refuses a write by a holder of the old epoch, atomically
with the effect, or only when the writer remembers to ask. One rule, applied to every mechanism from deterministic
probes, assigns an \emph{enforcement class}, strongest first:
\begin{itemize}
  \item \emph{surface-native}: a built-in check of the surface (a constraint, a permission) refuses the stale write, even
        when the writer omits every field that exists for the fence;
  \item \emph{author code in the surface}: code we installed and the surface runs on every write (a trigger) refuses it,
        under the same omission test;
  \item \emph{writer-supplied condition}: the surface refuses only because the writer submitted the condition (a compare,
        an epoch column); a writer that leaves it out commits;
  \item \emph{advisory}: the surface never refuses; writers that honour a lock refuse themselves.
\end{itemize}
Only the first two hold against a writer that is buggy or compromised, and the second only while the installed code is
part of the surface's trusted configuration.

\para{Probes (E6, E6b).} We probed local single-node reference instances of PostgreSQL~16, etcd~v3.5.17 and a POSIX
filesystem; Amazon S3 and the Kubernetes API are recorded from documentation only. E6 tested three primitives per
surface: idempotent retry (P1), an epoch fence atomic with the effect (P2), and a hold that survives an unclean stop and
that revocation cancels (P3). \vsixESixPOneThreeSentence{} Its P2 probe raced a stream of stale writers against the
fence and exposed a mechanism we had not anticipated: in PostgreSQL, new \texttt{FOR SHARE} lockers join an existing
share lock while the fence's \texttt{FOR UPDATE} waits, so a continuing stream of stale writers can starve the raise.

E6b, predeclared separately, re-tests only the fence. It probes
\vsixESixBMechanisms{} mechanisms deterministically, \vsixESixBRuns{} runs each, with \vsixESixBVariants{} writer
variants: the stale writer as specified, with the epoch
omitted, with every fence field omitted, claiming the new epoch, and a fresh writer. It then repeats E6's race
\vsixESixBStarvationN{} times per mechanism, with and without added load. \vsixESixBSentence{} Table~\ref{tab:v6fence}
gives every mechanism.

% Generated by scripts/v6_macros.py from evidence/v6/reference/e6b-fenceability; do not edit.
\begin{table*}[t]\centering\footnotesize
\caption{Epoch fences on effect surfaces (E6b; local single-node instances). \emph{Class} is assigned by one rule from deterministic probes: \emph{surface-native}, a built-in check refuses the stale write even when every fence field is omitted; \emph{author code}, code we installed in the surface does; \emph{writer}, the surface refuses only when the writer supplies the condition; \emph{advisory}, it never refuses. \emph{Fence} is the deterministic post-fence outcome; \emph{forged} counts stale writers claiming the new epoch that the surface refused (a dash where the fence could never be raised, so no new epoch existed to claim). The race columns repeat E6's race without added load: fence starved, at least one stale write refused by the surface, and no stale effect ordered after the fence. Amazon S3 and the Kubernetes API were not executed.}\label{tab:v6fence}
\begin{tabular}{lllllcccc}\toprule
Surface & Mechanism & Epoch in & Class & Fence & Forged & Starved & Refused & No stale after \\\midrule
PostgreSQL 16 & trigger (E6) & request & writer & pass & 0/40 & 18/20 & 6/20 & 20/20 \\
 & trigger + table lock & request & writer & pass & 0/40 & 0/20 & 20/20 & 20/20 \\
 & trigger, epoch NOT NULL & request & author code & pass & 0/40 & 16/20 & 9/20 & 20/20 \\
 & foreign key to live epoch & request & indeterminate & fail & -- & \multicolumn{3}{c}{fence raise failed (20/20)} \\
 & REVOKE writer role & credential & surface-native & pass & 40/40 & 0/20 & 20/20 & 0/20 \\
etcd v3.5.17 & txn compare on floor & request & writer & pass & 0/40 & 0/20 & 20/20 & 20/20 \\
 & RBAC revoke & credential & surface-native & pass & 40/40 & 0/20 & 20/20 & 20/20 \\
POSIX filesystem & flock + floor file & request & advisory & fail & 0/40 & 6/20 & 0/20 & 20/20 \\
 & directory mode bits & address & surface-native & pass & 0/40 & 0/20 & 20/20 & 20/20 \\
\bottomrule\end{tabular}\end{table*}


\para{Reading.} A fence needs three properties at once: the surface enforces it, a writer cannot claim the new epoch,
and it takes effect while stale writers are still running, with no stale write committing after it. A fence that is
acknowledged only once every writer has finished (starved, in Table~\ref{tab:v6fence}) has no write after it by
default, and we do not count it as holding. Of the \vsixFenceMechanisms{} mechanisms probed, only \vsixFenceAllThree{} had all
three (Figure~\ref{fig:v6hero}c), and the reason is specific to etcd. Each authenticated request carries the revision
of the authentication store at which its token was checked, any change to users, roles or permissions raises that
revision, and
the permission check for a write runs again when the write is applied, after consensus: a request whose revision is
older than the store's is refused there (\texttt{auth: revision in header is old})~\cite{etcdauth}. Revoking the
writer's role therefore reaches writes that were already accepted into the log, which a revocation at our broker
cannot do. The same check refuses legitimate requests issued before any unrelated change to authentication, which a
client retries with a fresh token. We probed one node, and \vsixFenceAllThreeRaces{} clean races bound the per-race
failure rate only below \vsixFenceRaceBound{} (one-sided 95\%); on a cluster each member applies the check on its own
apply path, which we did not test. Three results bear on the stop problem. First, fenceability is a property of the effect surface and of
how its fence is bound, not of the agent framework. An epoch that the writer states can be omitted or forged, so such a
fence stops a stale writer that is honest and nothing else. Second, binding the epoch to a credential resists a writer
that lies, but revoking a credential is not automatically a fence: on one surface it behaved exactly like revocation at
our broker, refusing new work while letting work already running commit. Whether a surface's revoke reaches in-flight
work has to be tested, not assumed. Third, a fence that is correct can still be slow to take effect under the load it is
meant to stop; its latency belongs in the stop budget of Section~\ref{sec:v6stop}. An operator who needs ``stop'' to
reach dispatched work must route effects through a surface whose fence is native and reaches in-flight writes, or place
a mediating sink in front of one whose fence does not.

\para{Limits of the census.} We chose the nine mechanisms by hand, one or a few per surface, and ran each on one
unloaded node; Amazon S3 and Kubernetes are classified from documentation only. The deterministic runs are repeated to
detect nondeterminism, not drawn as independent samples, so their counts are not rates. Refusing a forged epoch is
nearly definitional once the epoch is bound to a credential; the informative column is the third. We did not try
combinations, such as revoking a PostgreSQL role and terminating its sessions, which may close the gap that revocation
alone left; the trigger's class changed between E6 and E6b when the classification rule was fixed.

\subsection{What the specification does not name}

Section~\ref{sec:v6reviewers} showed that removing one clause of a gate admits exactly the violations that break only
that clause, and Section~\ref{sec:v6cost} measured what injections that stay inside a written specification achieve.
Neither is a defect that more verification would remove. The kernel's proofs establish that the
implementation enforces its predicate; they cannot establish that the predicate is the one the operator meant. What
mediation buys is relocation of the error, from a learned judgement that an adversary probes by search to a written
artefact that a person can read, diff, test against adversarial suites, and amend. We regard that relocation, not the
zeros, as the contribution.

\subsection{Intelligence re-enters through conferral}

Everything above holds a grant fixed and asks what a proposer can do within it. Grants are not fixed. People widen
them for agents that appear competent and careful. That creates a route by which capability can turn back into
authority, and nothing in this paper's mechanism blocks it.

\begin{observation}[Conferral reversal]\label{obs:conferral}
Let \(h_d(g)\) be the expected harm of disposition \(d\) under grant \(g\), and \(\pi_d\) the distribution of grants
that operators confer on \(d\). It is possible that \(h_A(g) < h_B(g)\) for every \(g\) and yet
\(\mathbb{E}_{g\sim\pi_A} h_A(g) > \mathbb{E}_{g\sim\pi_B} h_B(g)\), whenever \(A\) earns sufficiently broader grants.
\end{observation}

This is the mixture reversal familiar from Simpson's paradox, applied to authority. It is not a result about any model.
Whether it occurs depends on how real operators confer authority, which requires a study with people. We predeclare its
design and cost in Appendix~\ref{app:v6conferral}: operators grant budgets to agents of controlled competence and
controlled caution, with harm measured at a mediated sink. We do not run it in this version. Until it is run, the claim
that mediation makes reach independent of the proposer must be read as ``at a fixed grant''.
